% SPDX-License-Identifier: CC-BY-4.0
% Copyright (c) 2026 Martin Schröder <info@swedishembedded.com>

\section{Introduction}

A team that wants a local coding model to be competent in its own codebase has two
unsatisfying options: prompt a general model with the repository, or fine-tune on data
somebody curated by hand. Recent work shows a third way is possible for Python
projects: derive executable tasks from a repository, let an agent attempt them, keep the
attempts that hidden tests confirm, and train on those
\citep{yang2025swesmith,pan2025swegym,jain2025r2egym}. That work trains on pooled
corpora spanning many repositories and evaluates on SWE-bench-style benchmarks
\citep{jimenez2024swebench}. Our question is the operational one that follows.

\paragraph{Research question.} Can we take any git repository and tune a local agentic
model to work competently in the domain of that repository; can we do this repeatedly,
with further repositories and with later revisions of the same one, at any point in a
project's life; and does the model get progressively better in the domains it has
met, rather than better on the newest one and worse on the rest?

\paragraph{Thesis.} Under an identical agent harness, a LoRA adapter \citep{hu2021lora}
trained only on hidden-test-verified trajectories of tasks derived from repository $R$
resolves more of $R$'s reserved task families than the base model does; and after an
adapter has further been trained on other repositories, it resolves no fewer of
earlier repositories' reserved families than it did after their own round, within a
margin fixed in advance.

\paragraph{What would count against it.} (i) No significant paired gain on $R$'s reserved
families, or a loss. (ii) A gain that a control explains without learning the
repository: an adapter trained on a \emph{different} repository gains as much on $R$
(it learned the tool protocol, not $R$), or an adapter trained on unverified
trajectories gains as much (verification did not matter). (iii) A drop on earlier
repositories beyond the margin after later rounds.

\paragraph{Contributions.}
\begin{enumerate}
\item A design in which software engineering is a general \emph{workspace} task kind of
a learning system whose other task kinds (documents, tools, persons) share its
experience store, curriculum, release gate and lineage, rather than a separate
benchmark pipeline (\cref{sec:system}).
\item A frozen measurement protocol for the repeated, multi-repository case: families
reserved before any rollout, a cluster-paired exact sign test as the primary decision,
retention as a non-inferiority check, cross-repository and unverified-data controls, and an exposure
audit between what the solver sees and what the verifier holds (\cref{sec:protocol}),
with a table of what that test can decide at the sample sizes one repository yields.
\item The defects found by auditing the existing stack against that protocol
(\cref{sec:audit}).
\end{enumerate}

\paragraph{Status.} This version contains no learning result. The protocol was fixed
before any measured run; every experiment in \cref{sec:results} is \notrun\ until its row
exists in \code{data/results.csv}.
